% Generated draft: requires mathematical and source review.
\documentclass[11pt]{article}
\usepackage[margin=25mm]{geometry}
\usepackage{fontspec}
\setmainfont{Latin Modern Roman}
\usepackage{amsmath,amssymb,mathtools,unicode-math}
\setmathfont{Latin Modern Math}
\usepackage{enumitem,needspace,xurl,hyperref}
\setlength{\parindent}{0pt}
\setlength{\parskip}{5pt}
\setlength{\emergencystretch}{3em}
\newcommand{\N}{\mathbb{N}}
\newcommand{\Z}{\mathbb{Z}}
\newcommand{\Q}{\mathbb{Q}}
\newcommand{\R}{\mathbb{R}}
\newcommand{\C}{\mathbb{C}}
\newcommand{\F}{\mathbb{F}}
\newcommand{\A}{\mathbb{A}}
\newcommand{\PP}{\mathbb{P}}
\newcommand{\E}{\mathbb{E}}
\newcommand{\eps}{\varepsilon}
\newcommand{\dd}{\,\mathrm{d}}
\begin{document}
\section*{UnsolvedMath ID: 30002887}
\textbf{Source problem code:} catalogue entry
\subsection{Definitions and mathematical statement}
A conventional framework for interpreting the request is as follows; the source does not specify all hypotheses needed to adopt it. Over an algebraically closed field \(k\), a symmetric homogeneous space is usually a quotient \(G/H\), where \(G\) is a connected reductive algebraic group, \(\theta\colon G\to G\) is an involutive automorphism, and \(H\) is a closed subgroup satisfying
\[
(G^\theta)^\circ\subseteq H\subseteq G^\theta,
\qquad G^\theta=\{g\in G:\theta(g)=g\}.
\]
The source explicitly requires \(G\) to be reductive, but does not specify the base field, its characteristic, connectedness, the involution, or the precise permitted subgroups \(H\). The displayed convention is therefore explanatory rather than an additional source hypothesis.

An equivariant compactification of \(G/H\) is a complete normal \(G\)-variety \(X\) with a \(G\)-equivariant open immersion \(G/H\hookrightarrow X\) having dense image. In the spherical setting, choose a Borel subgroup \(B\subseteq G\). A color is a \(B\)-stable prime divisor of \(G/H\) that is not \(G\)-stable. The embedding is toroidal if the closure in \(X\) of every color contains no \(G\)-orbit. These definitions explain the requested toroidal compactifications; the source does not supply smoothness or projectivity assumptions.

Over an algebraically closed field, a rational curve is an integral curve whose normalization is isomorphic to \(\mathbb P^1_k\). The expression \emph{minimal rational curves} requires a specified convention concerning families of rational curves and their minimality. The source does not specify such a convention, a polarization, or a numerical degree criterion.

\textbf{Requested extension.} Develop a theory of minimal rational curves on complete symmetric varieties, specifically on toroidal compactifications \(X\) of symmetric spaces \(G/H\) with \(G\) reductive, extending the results for wonderful group compactifications.

Here the group case conventionally concerns a reductive group \(K\) as the homogeneous space \((K\times K)/\operatorname{diag}(K)\), with left and right multiplication. The source does not identify the particular wonderful group compactifications, hypotheses on \(K\), or results to be extended. Consequently, the supplied request is a research program rather than a uniquely specified conjecture: its exact conclusions and quantifiers cannot be recovered from this record. In particular, the phrase ``complete symmetric varieties'' must retain the source's stated restriction to toroidal compactifications; completeness alone does not express the toroidal condition.
\subsection{Short English statement}
The problem asks for an extension of the theory of minimal rational curves from wonderful group compactifications to complete toroidal compactifications of reductive symmetric spaces \(G/H\). The field, geometric regularity assumptions, meaning of minimality, and particular results to extend remain unspecified, so a single precise theorem cannot yet be stated. The supplied record labels the problem open; this draft does not independently verify that status.
\subsection{Sources}
\begin{itemize}
\item Local UnsolvedMath registry, \texttt{data/unsolvedmath/problems.json}, permanent ID \(30002887\), title ``Minimal Rational Curves on Complete Symmetric Varieties,'' exact source problem code \texttt{catalogue entry}. The supplied record gives proposed year \(2015\) and names no proposer.
\item Supplied report citation: \emph{Enveloping Algebras and Geometric Representation Theory} (2016), Oberwolfach Reports, pp.\ 1385--1447, DOI \texttt{10.4171/owr/2015/25}. \url{https://doi.org/10.4171/owr/2015/25}. The background attributes the problem to workshop open-problem discussions recorded in Oberwolfach Reports.
\item Supplied literature source: Ruzzi, \emph{Fano symmetric varieties with low rank}. \url{https://ems.press/content/serial-article-files/41168}.
\item Supplied literature source: \emph{Oberwolfach Report 2015/25}. \url{https://publications.mfo.de/bitstream/handle/mfo/3472/OWR_2015_25.pdf?isAllowed=y&sequence=1}.
\item Literature assessment dated exactly \texttt{2026-08-21}: the supplied curation labels the problem open and reports a substantial colored-fan and low-rank theory for complete and toroidal symmetric varieties, while describing the requested general extension to all symmetric spaces \(G/H\) as open. This is supplied dated curation derived from the OWR record, not an independently verified literature assessment; its time-sensitive status and cited claims require independent verification.
\item The registry marks the statement \texttt{corrected\_verified} and reports a manual semantic, typographic, and LaTeX review of its cleanup. This supplied editorial qualification does not establish mathematical correctness or identify the precise results to be generalized.
\end{itemize}
\end{document}
